\section{Fault model}
\label{sec:faults}

\subsection{Fault types}

A fault is a deviation applied to one tool call.
Pre-execution faults stop the call before it runs; post-execution faults change the result of a call that already ran, so its side effects have happened.
Table~\ref{tab:faults} lists the types used in the headline profile.

\begin{table}[t]
\centering
\small
\begin{tabular}{llp{7.2cm}l}
\toprule
Type & Phase & What the agent sees & Layers \\
\midrule
\code{deny} & pre & Call refused, permission-style message & S, N \\
\code{err} & pre & Timeout, unavailable, connection reset & S, N \\
\code{rate} & pre & HTTP 429 with \code{Retry-After} & S \\
\code{empty} & post & Result replaced by an empty string & S, N \\
\code{trunc} & post & Result cut to its first 30\% with no marker & S, N \\
\code{falsefail} & post & Call executed, result reports an error & S, N \\
\code{poison} & post & Plausible stale data carrying a unique token & S, N \\
\code{drift} & post & JSON shape changed (renamed key, nesting, number as string) & S \\
\bottomrule
\end{tabular}
\caption{Fault types in the \code{moderate} profile. S is the service track, N the native track.}
\label{tab:faults}
\end{table}

The specification defines three more types that the headline profile does not use: input corruption on read, search and edit arguments, latency, and silent write corruption.
\code{falsefail} applies on the service track only to tools the task declares non-idempotent, so it tests whether the agent repeats a side effect.
\code{poison} applies only to resources the task declares, with an author-written stale variant; on shell calls only the lines the variant changes are replaced, where the output shows them, and the command itself is never changed.

Payload text mimics real tooling (errno codes, HTTP status lines, git and npm style errors) and comes from one versioned catalog.
A lint rule keeps words that would reveal the experiment out of the catalog, injector source, file names and environment variable names.

\subsection{Schedule}

Whether a call is faulted is a pure function of the call, not of timing or order.
For every eligible call,
\begin{align*}
k &= \mathrm{sha256}(\mathit{seed} \,\|\, \mathit{task} \,\|\, \mathit{attempt} \,\|\, \mathit{class} \,\|\, \mathit{args}),\\
u &= \mathrm{uint32}(\mathrm{sha256}(k \,\|\, \mathit{occurrence})) / 2^{32},
\end{align*}
and the call is faulted when $u$ is below the profile's rate for its capability class, the per-trial cap is not reached, and the call is not exempt.
The fault type is a weighted choice from a second hash, restricted to the harness's parity set; the payload is a deterministic choice among matching catalog entries.
Arguments are canonicalized first: paths relative to the workspace, JSON with sorted keys, shell commands with collapsed whitespace.
Because attempt $a$ of every configuration draws from the same function, two harnesses that issue the same call in the same attempt receive the same decision.
On the service track tool names and arguments are shared, so identical calls are common; on the native track parity is distributional (same rates, type mix, persistence and caps).

\subsection{Persistence and targeted faults}

A transient fault fails the first occurrence of a key; \code{sticky($k$)} fails the first $k$ consecutive occurrences.
Separately from the per-call draw, a targeted draw selects, once per trial, whether each declared poisonable resource is stale for the whole trial and, on the service track, whether the first call to each declared non-idempotent tool reports a false failure.
Faults never touch the agent's status file, the verifier, the oracle, logs, harness configuration or injector directories.

\subsection{Solvability}

Bounded transient and sticky faults leave every task solvable by an agent that retries or reroutes.
On the service track every oracle solution calls the services through the fault proxy with bounded retries and idempotency keys, and must pass under every profile before a task is accepted.
A resource that stays stale for the whole trial always leaves a route to the true data: on the native track the file on disk is unchanged, and on the service track another call returns the true record.
